\begin{figure}[htbp]
\centering
\begin{tikzpicture}
\begin{axis}[
    width=13cm, height=6.5cm,
    xlabel={Frequency (rad/s)},
    ylabel={Magnitude (dB)},
    xmin=0.1, xmax=10000,
    ymin=-80, ymax=5,
    xmode=log,
    grid=major,
    legend pos=south left,
    legend style={font=\footnotesize},
    title={Filter Type Comparison: Magnitude Response (Normalized Cutoff $\omega_c = 100$ rad/s)},
    title style={font=\small, color=UofSCBlack}
]

% Low-pass filter
\addplot[UofSCGarnet, thick, domain=0.1:10000, samples=200] {20*log10(1/sqrt(1 + (x/100)^2))};
\addlegendentry{Low-Pass}

% High-pass filter
\addplot[UofSCAtlantic, thick, domain=0.1:10000, samples=200] {20*log10((x/100)/sqrt(1 + (x/100)^2))};
\addlegendentry{High-Pass}

% Band-pass filter (Q=5, centered at 100 rad/s)
\addplot[UofSCHorseshoe, thick, domain=0.1:10000, samples=200] {20*log10((x/100)/sqrt((1-(x/100)^2)^2 + (x/(100*5))^2))};
\addlegendentry{Band-Pass}

% Cutoff frequency marker
\draw[dashed, UofSC50Black, thin] (axis cs:100, -80) -- (axis cs:100, 5);
\node[anchor=north, color=UofSC50Black, font=\footnotesize] at (axis cs:100, -80) {$\omega_c$};

\end{axis}
\end{tikzpicture}
\caption{Comparison of filter types showing magnitude response of low-pass, high-pass, and band-pass filters with normalized cutoff frequency.}
\label{fig:filter_comparison}
\end{figure}
